% C API reference.

\guidechapter{C API REFERENCE}{%
  \item Every public structure, function and capability macro
  \item Exact prototypes, parameters, return values and side effects
  \item Short, compilable call examples
}

\begin{referencechapter}

\newcommand{\capiref}[8]{%
  \capi{#1}%
  \entryfield{Header}{\cmd{#2}}%
  \entryfield{Prototype}{\cmd{#3}}%
  \entryhead{Parameters}#4\par
  \entryhead{Returns}#5\par
  \entryhead{Purpose}#6\par
  \entryhead{Example}%
  \begin{listingbox}[before skip=0.3\baselineskip,after skip=0.3\baselineskip]#7\end{listingbox}%
  \begin{resultbox}#8\end{resultbox}}

\newcommand{\cmacroref}[7]{%
  \cmacroapi{#1}%
  \entryfield{Header}{\cmd{ultimate.h}}%
  \entryfield{Syntax}{\cmd{#2}}%
  \entryhead{Parameters}#3\par
  \entryhead{Returns}#4\par
  \entryhead{Purpose}#5\par
  \entryhead{Example}%
  \begin{listingbox}[before skip=0.3\baselineskip,after skip=0.3\baselineskip]#6\end{listingbox}%
  \begin{resultbox}#7\end{resultbox}}

\newcommand{\capirefout}[8]{%
  \capi{#1}%
  \entryfield{Header}{\cmd{#2}}%
  \entryfield{Prototype}{\cmd{#3}}%
  \entryhead{Parameters}#4\par
  \entryhead{Returns}#5\par
  \entryhead{Purpose}#6\par
  \entryhead{Example}%
  \begin{listingbox}[before skip=0.3\baselineskip,after skip=0.3\baselineskip]#7\end{listingbox}%
  \begin{outputbox}#8\end{outputbox}}

\newcommand{\structref}[3]{%
  \apientry{#1 type}{\cmd{#1}}{C structure}%
  \label{struct:#2}%
  \entryfield{Header}{\cmd{#3}}}

\section{Common Rules}

Include \cmd{ultimate.h} for application services and \cmd{uci.h} only for raw
transport calls. Unless an entry says otherwise, a \cmd{uint8\_t} service
returns \cmd{ULTIMATE\_OK} on success or an \cmd{ULTIMATE\_ERR\_*} value from
Appendix~A. Caller-supplied buffers remain owned by the caller and must stay
valid until the call returns. The SDK allocates no memory.
For UCI-backed calls, an entry's Returns field names call-specific outcomes;
transport and target failures may also return the applicable Appendix~A error.

Call \cmd{ultimate\_init()} once before any UCI-backed service. Turbo, REU DMA
and software vsprites state their separate requirements. cc65 character maps
string literals; uppercase filename/path literals are ASCII-compatible, while
case-sensitive ASCII URLs, headers, keys and values should be constructed as
bytes or converted at run time.

\section{SDK Data Types}

The following entries define every public structure used by the C API. Function
parameter descriptions link back here; the structure definitions in
\cmd{ultimate.h} and \cmd{uci.h} remain the compiler's authority.

\structref{ultimate\_capabilities}{ultimate-capabilities}{ultimate.h}

Filled by \cmd{ultimate\_detect()}; callers normally allocate it once and keep
it for feature checks.

\begin{apitable}{p{0.25\linewidth}p{0.67\linewidth}}
\textbf{Field} & \textbf{Meaning} \\ \midrule
\cmd{uint8\_t present} & 1 when a working command interface answered. \\
\cmd{uint8\_t ident} & Raw UCI identification-register value. \\
\cmd{uint16\_t targets} & Bit N is set when UCI target N identified itself. \\
\end{apitable}

\entryhead{Example}
\begin{listingbox}
ultimate\_capabilities caps;\
uint8\_t err = ultimate\_detect(\&caps);\
if (err == ULTIMATE\_OK \&\& ultimate\_has\_http(\&caps)) \{\
\hspace*{2em}/* HTTP target is available */\
\}
\end{listingbox}

\structref{ultimate\_sid}{ultimate-sid}{ultimate.h}

One raw SID mapping returned inside \hyperref[struct:ultimate-sid-info]{%
\cmd{ultimate\_sid\_info}}. It describes addresses and firmware mapping type;
it does not identify a 6581 versus 8580.

\begin{apitable}{p{0.34\linewidth}p{0.58\linewidth}}
\textbf{Field} & \textbf{Meaning} \\ \midrule
\cmd{uint16\_t primary\_address} & Primary memory-mapped SID address. \\
\cmd{uint16\_t secondary\_address} & Second mapping, or zero when absent. \\
\cmd{uint8\_t type} & Raw implementation/address-mask code; preserve unknown values. \\
\end{apitable}

\entryhead{Example}
\begin{listingbox}
ultimate\_sid\_info info;\
ultimate\_sid *sid;\
uint8\_t err = ultimate\_legacy\_get\_sid\_info(\&info);\
if (err == ULTIMATE\_OK \&\& info.count != 0)\
\hspace*{2em}sid = \&info.sid[0];
\end{listingbox}

\structref{ultimate\_sid\_info}{ultimate-sid-info}{ultimate.h}

Count-prefixed result from deprecated control HWINFO. At most
\cmd{ULTIMATE\_SID\_MAX} (4) records are retained.

\begin{apitable}{p{0.34\linewidth}p{0.58\linewidth}}
\textbf{Field} & \textbf{Meaning} \\ \midrule
\cmd{uint8\_t count} & Valid entries in \cmd{sid}; zero on any call failure. \\
\cmd{ultimate\_sid sid[ULTIMATE\_SID\_MAX]} & SID mapping records; only
indices below \cmd{count} are valid. \\
\end{apitable}

\entryhead{Example}
\begin{listingbox}
ultimate\_sid\_info info;\
uint8\_t err = ultimate\_legacy\_get\_sid\_info(\&info);\
if (err == ULTIMATE\_OK \&\& info.count != 0) \{\
\hspace*{2em}/* info.sid[0] is valid */\
\}
\end{listingbox}

\structref{ultimate\_vsprite}{ultimate-vsprite}{ultimate.h}

A complete software-vsprite draw descriptor. It is read only during
\cmd{ultimate\_vsprite\_draw()}; the bitmap, and optionally screen memory, are
changed in place. The descriptor is 14 bytes in the cc65 ABI.

\begin{apitable}{p{0.27\linewidth}p{0.65\linewidth}}
\textbf{Field} & \textbf{Meaning} \\ \midrule
\cmd{uint8\_t *bitmap} & Base of a standard 8000-byte C64 bitmap. \\
\cmd{const uint8\_t *source} & Image bytes in column-major order; for
\cmd{VSPRITE\_F\_COPY}, the base of another C64 bitmap. \\
\cmd{const uint8\_t *mask} & Column-major mask required by
\cmd{VSPRITE\_F\_MASKED} and ignored otherwise; the destination is ANDed with
the mask, then the source is ORed into it. \\
\cmd{uint8\_t *screen} & Screen-memory base required by
\cmd{VSPRITE\_F\_COLOR} and ignored otherwise. \\
\cmd{uint8\_t x} & Destination bitmap byte/cell column, 0--39; this is not a
pixel coordinate. \\
\cmd{uint8\_t y} & Destination raster line, 0--199. \\
\cmd{uint8\_t width} & Width in bitmap byte columns; must be nonzero and fit
from \cmd{x}. \\
\cmd{uint8\_t height} & Height in raster lines; must be nonzero and fit from
\cmd{y}. \\
\cmd{uint8\_t color} & Screen byte written to cells touched by visible source
bytes in colour mode. \\
\cmd{uint8\_t flags} & 0 for OR, or exactly one of
\cmd{VSPRITE\_F\_MASKED}, \cmd{VSPRITE\_F\_COPY}, or
\cmd{VSPRITE\_F\_COLOR}. \\
\end{apitable}

For ordinary and masked draws, source (and mask) layout is every raster line of
byte column zero, followed by every line of byte column one, and so on. The
caller supplies pre-shifted images when sub-cell horizontal movement is needed.

\Needspace{22\baselineskip}
\entryhead{Complete masked-draw example}
\begin{listingbox}
static const uint8\_t image[8] =\
\hspace*{2em}\{ 0x18, 0x3c, 0x7e, 0xff, 0xff, 0x7e, 0x3c, 0x18 \};\
static const uint8\_t mask[8] =\
\hspace*{2em}\{ 0xe7, 0xc3, 0x81, 0x00, 0x00, 0x81, 0xc3, 0xe7 \};\
ultimate\_vsprite sprite;\
uint8\_t err;\
\
sprite.bitmap = (uint8\_t *)0xe000;\
sprite.source = image;\
sprite.mask = mask;\
sprite.screen = (uint8\_t *)0xcc00;\
sprite.x = 18;       /* byte column */\
sprite.y = 80;       /* raster line */\
sprite.width = 1;    /* 8 pixels */\
sprite.height = 8;\
sprite.color = 2;    /* ignored in masked mode */\
sprite.flags = VSPRITE\_F\_MASKED;\
err = ultimate\_vsprite\_draw(\&sprite);
\end{listingbox}

\begin{resultbox}
The 8-by-8 image is composited at bitmap byte column 18, raster line 80. Set
\cmd{flags} to 0, \cmd{VSPRITE\_F\_COPY}, or \cmd{VSPRITE\_F\_COLOR} for the
other three operations and populate the fields described above.
\end{resultbox}

\structref{ultimate\_fileinfo}{ultimate-fileinfo}{ultimate.h}

Metadata returned by \cmd{ultimate\_stat()} and \cmd{ultimate\_fstat()}.

\begin{apitable}{p{0.31\linewidth}p{0.61\linewidth}}
\textbf{Field} & \textbf{Meaning} \\ \midrule
\cmd{uint32\_t size} & File size in bytes. \\
\cmd{uint16\_t date} & Packed FAT date: year minus 1980, month and day. \\
\cmd{uint16\_t time} & Packed FAT time: hour, minute and two-second units. \\
\cmd{char extension[3]} & Three space-padded bytes; not NUL-terminated. \\
\cmd{uint8\_t attrib} & \cmd{DOS\_ATTR\_*} bits. \\
\cmd{char name[ULTIMATE\_NAME\_MAX + 1]} & NUL-terminated name, bounded by
\cmd{ULTIMATE\_NAME\_MAX}. \\
\end{apitable}

\entryhead{Example}
\begin{listingbox}
ultimate\_fileinfo info;\
uint8\_t err = ultimate\_stat("DEMO.PRG", \&info);\
if (err == ULTIMATE\_OK) \{\
\hspace*{2em}/* info.name, info.size and info.attrib are valid */\
\}
\end{listingbox}

\structref{ultimate\_audio\_voice}{ultimate-audio-voice}{ultimate.h}

One Ultimate Audio PCM channel descriptor. Configuration reads all fields and
leaves the channel stopped. The descriptor is \cmd{UA\_VOICE\_SIZE} (22)
bytes in the cc65 ABI.

\begin{apitable}{p{0.29\linewidth}p{0.63\linewidth}}
\textbf{Field} & \textbf{Meaning} \\ \midrule
\cmd{uint8\_t channel} & Channel 0 through \cmd{UA\_CHANNELS-1}. \\
\cmd{uint8\_t flags} & \cmd{UA\_CTRL\_*} format/repeat/IRQ bits, excluding
\cmd{UA\_CTRL\_GATE}. \\
\cmd{uint8\_t volume} & 0 through \cmd{UA\_VOLUME\_MAX}. \\
\cmd{uint8\_t pan} & \cmd{UA\_PAN\_LEFT} through \cmd{UA\_PAN\_RIGHT}. \\
\cmd{uint32\_t reu\_address} & First PCM byte in the 24-bit REU window. \\
\cmd{uint32\_t length} & PCM byte count, 1 through \$FFFFFF. \\
\cmd{uint32\_t repeat\_a} & Loop-start byte offset relative to
\cmd{reu\_address}. \\
\cmd{uint32\_t repeat\_b} & Loop-end byte offset relative to
\cmd{reu\_address}. \\
\cmd{uint16\_t rate} & Nonzero divider: round(6250000 / samples per second). \\
\end{apitable}

Values 7 and 8 are both centred pan positions. The address and byte count may
end at 16\,MB but may not wrap beyond it. A 16-bit or interleaved voice needs
an even length. With \cmd{UA\_CTRL\_REPEAT}, the required relation is
\cmd{repeat\_a < repeat\_b <= length}.

\Needspace{18\baselineskip}
\entryhead{Complete 8-bit mono example}
\begin{listingbox}
ultimate\_audio\_voice voice;\
uint8\_t err;\
voice.channel = 0;\
voice.flags = 0;\
voice.volume = UA\_VOLUME\_MAX;\
voice.pan = UA\_PAN\_CENTER;\
voice.reu\_address = 0x4000UL;\
voice.length = 8000UL;\
voice.repeat\_a = 0;\
voice.repeat\_b = 0;\
voice.rate = 781;       /* about 8000 samples/second */\
err = ultimate\_audio\_configure(\&voice);\
if (err == ULTIMATE\_OK)\
\hspace*{2em}err = ultimate\_audio\_start(voice.channel, voice.flags);
\end{listingbox}

\structref{ultimate\_drive}{ultimate-drive}{ultimate.h}

One drive or occupied IEC-bus record inside
\hyperref[struct:ultimate-drives]{\cmd{ultimate\_drives}}.

\begin{apitable}{p{0.26\linewidth}p{0.66\linewidth}}
\textbf{Field} & \textbf{Meaning} \\ \midrule
\cmd{uint8\_t type} & \cmd{CTRL\_DRVTYPE\_*}; SoftwareIEC and printer values
describe bus slots, not physical drives. \\
\cmd{uint8\_t device} & IEC device number accepted by mount operations. \\
\cmd{uint8\_t power} & 1 while the drive is running, otherwise 0. \\
\end{apitable}

\entryhead{Example}
\begin{listingbox}
ultimate\_drives drives;\
ultimate\_drive *drive;\
uint8\_t err = ultimate\_drive\_info(\&drives);\
if (err == ULTIMATE\_OK \&\& drives.count != 0)\
\hspace*{2em}drive = \&drives.drive[0];
\end{listingbox}

\structref{ultimate\_drives}{ultimate-drives}{ultimate.h}

Count-prefixed result from \cmd{ultimate\_drive\_info()}.

\begin{apitable}{p{0.38\linewidth}p{0.54\linewidth}}
\textbf{Field} & \textbf{Meaning} \\ \midrule
\cmd{uint8\_t count} & Valid records; never greater than records actually
received or \cmd{ULTIMATE\_DRIVES\_MAX}. \\
\cmd{ultimate\_drive drive[ULTIMATE\_DRIVES\_MAX]} & Drive/bus records; only indices below
\cmd{count} are valid. \\
\end{apitable}

\entryhead{Example}
\begin{listingbox}
ultimate\_drives drives;\
uint8\_t err = ultimate\_drive\_info(\&drives);\
if (err == ULTIMATE\_OK \&\& drives.count != 0) \{\
\hspace*{2em}/* drives.drive[0] is valid */\
\}
\end{listingbox}

\structref{uci\_request}{uci-request}{uci.h}

One low-level command/reply exchange. All pointed-to spans remain owned by the
caller. The cc65 layout is \cmd{UCI\_REQ\_SIZE} (22) bytes and matches the
generated assembly offsets.

\begin{apitable}{p{0.27\linewidth}p{0.65\linewidth}}
\textbf{Field} & \textbf{Meaning} \\ \midrule
\cmd{uint8\_t target} & \cmd{UCI\_TARGET\_*}, optionally ORed with
\cmd{UCI\_TARGET\_NO\_REPLY}. \\
\cmd{uint8\_t command} & Target-specific command byte. \\
\cmd{const uint8\_t *args} & Optional fixed argument bytes. \\
\cmd{uint16\_t arglen} & Byte count at \cmd{args}; use zero with NULL. \\
\cmd{const uint8\_t *payload} & Optional payload span sent after arguments. \\
\cmd{uint16\_t payloadlen} & Byte count at \cmd{payload}; use zero with NULL. \\
\cmd{uint8\_t *data} & Reply-data buffer, or NULL to discard without a
truncation result. \\
\cmd{uint16\_t datamax} & Capacity of \cmd{data}. \\
\cmd{uint16\_t datalen} & Output: reply-data bytes retained. \\
\cmd{uint8\_t *status} & Optional raw status buffer. \\
\cmd{uint16\_t statusmax} & Capacity of \cmd{status}. \\
\cmd{uint16\_t statuslen} & Output: status bytes retained. \\
\end{apitable}

\Needspace{17\baselineskip}
\entryhead{Complete IDENTIFY example}
\begin{listingbox}
uint8\_t reply[64];\
uint8\_t status[16];\
uci\_request req = \{ 0 \};\
uint8\_t err;\
req.target = UCI\_TARGET\_DOS1;\
req.command = UCI\_CMD\_IDENTIFY;\
req.data = reply;\
req.datamax = sizeof reply;\
req.status = status;\
req.statusmax = sizeof status;\
err = uci\_exec(\&req);
\end{listingbox}

\begin{resultbox}
On success, \cmd{reply[0..req.datalen-1]} contains the identification data and
\cmd{status[0..req.statuslen-1]} contains the raw target status.
\end{resultbox}

\section{Initialization And Discovery}

\capiref{ultimate\_init}{ultimate.h}
{uint8\_t ultimate\_init(void)}
{None.}
{\cmd{ULTIMATE\_OK}, \cmd{ULTIMATE\_ERR\_NO\_DEVICE}, or
\cmd{ULTIMATE\_ERR\_TIMEOUT}.}
{Initializes the SDK and functionally checks the command interface. Call it
once before other UCI-backed services.}
{uint8\_t err = ultimate\_init();}
{The SDK is ready only when \cmd{err == ULTIMATE\_OK}.}

\capiref{ultimate\_available}{ultimate.h}
{uint8\_t ultimate\_available(void)}
{None.}
{1 after successful initialization; otherwise 0.}
{Reads cached SDK state without touching hardware.}
{if (!ultimate\_available()) return 1;}
{The program exits when initialization has not succeeded.}

\capiref{ultimate\_detect}{ultimate.h}
{uint8\_t ultimate\_detect(ultimate\_capabilities *caps)}
{\cmd{caps} points to writable
\hyperref[struct:ultimate-capabilities]{\cmd{ultimate\_capabilities}} storage.}
{\cmd{ULTIMATE\_OK} after probing all known targets, including when some are
absent; \cmd{ULTIMATE\_ERR\_INVALID\_ARGUMENT} for a null \cmd{caps}; or the
\cmd{ULTIMATE\_ERR\_NO\_DEVICE}/\cmd{ULTIMATE\_ERR\_TIMEOUT} result from
initializing the transport.}
{Fills \cmd{present}, \cmd{ident}, and the target bit mask. It costs one UCI
round trip per known target, so normally call it once.}
{ultimate\_capabilities caps; err = ultimate\_detect(\&caps);}
{\cmd{caps.targets} identifies the services actually implemented.}

\capiref{ultimate\_identify}{ultimate.h}
{uint8\_t ultimate\_identify(uint8\_t target, char *buf, uint16\_t buflen, uint16\_t *outlen)}
{\cmd{target} is \cmd{UCI\_TARGET\_*}; \cmd{buf} has \cmd{buflen} bytes;
\cmd{outlen} may be \cmd{NULL}.}
{\cmd{ULTIMATE\_OK}, or \cmd{ULTIMATE\_ERR\_NOT\_SUPPORTED} when the target is
absent.}
{Reads a target's identification string, stores at most \cmd{buflen-1} bytes,
and terminates it. A supplied \cmd{outlen} excludes the terminator.}
{err = ultimate\_identify(UCI\_TARGET\_DOS1, text, sizeof text, \&length);}
{\cmd{text} contains the DOS identity and \cmd{length} its stored length.}

\capiref{ultimate\_get\_model}{ultimate.h}
{uint8\_t ultimate\_get\_model(char *buf, uint16\_t buflen, uint16\_t *outlen)}
{\cmd{buf} has \cmd{buflen} bytes; \cmd{outlen} may be \cmd{NULL}.}
{\cmd{ULTIMATE\_OK} or \cmd{ULTIMATE\_ERR\_NOT\_SUPPORTED}.}
{Reads the hardware model through the deprecated control HWINFO command.
Retained for compatibility; new code must handle its disappearance.}
{err = ultimate\_get\_model(text, sizeof text, NULL);}
{\cmd{text} is terminated when the call succeeds.}

\capiref{ultimate\_legacy\_get\_sid\_info}{ultimate.h}
{uint8\_t ultimate\_legacy\_get\_sid\_info(ultimate\_sid\_info *info)}
{\cmd{info} points to writable
\hyperref[struct:ultimate-sid-info]{\cmd{ultimate\_sid\_info}} storage.}
{\cmd{ULTIMATE\_OK}, \cmd{ULTIMATE\_ERR\_NOT\_SUPPORTED},
\cmd{ULTIMATE\_ERR\_PROTOCOL}, or \cmd{ULTIMATE\_ERR\_TRUNCATED};
\cmd{info->count} is zero on failure.}
{Reads up to four SID mapping records through deprecated HWINFO. Type bytes are
raw implementation/mapping codes, not 6581/8580 model identifiers.}
{ultimate\_sid\_info info; err = ultimate\_legacy\_get\_sid\_info(\&info);}
{On success, \cmd{info.sid[0].primary\_address} is the first SID address.}

\section{Capability Macros}

\cmacroref{ultimate\_has\_target}
{ultimate\_has\_target(caps, target)}
{\cmd{caps} points to
\hyperref[struct:ultimate-capabilities]{\cmd{ultimate\_capabilities}} filled by
\cmd{ultimate\_detect()}; \cmd{target} is a target number.}
{1 when that target bit is set, otherwise 0.}
{Tests an arbitrary target without another hardware transaction.}
{if (ultimate\_has\_target(\&caps, UCI\_TARGET\_HTTP)) use\_http();}
{The branch runs only when detection found the HTTP target.}

\cmacroref{ultimate\_has\_dos}{ultimate\_has\_dos(caps)}
{\cmd{caps} points to detected
\hyperref[struct:ultimate-capabilities]{\cmd{ultimate\_capabilities}}.}
{1 when DOS target 1 exists.}
{Tests the primary filesystem service.}
{if (ultimate\_has\_dos(\&caps)) list\_files();}
{Filesystem code runs only when DOS1 was detected.}

\cmacroref{ultimate\_has\_dos2}{ultimate\_has\_dos2(caps)}
{\cmd{caps} points to detected
\hyperref[struct:ultimate-capabilities]{\cmd{ultimate\_capabilities}}.}
{1 when DOS target 2 exists.}
{Tests the optional second filesystem service.}
{if (ultimate\_has\_dos2(\&caps)) inspect\_second\_dos();}
{The branch is skipped on machines exposing only DOS1.}

\cmacroref{ultimate\_has\_network}{ultimate\_has\_network(caps)}
{\cmd{caps} points to detected
\hyperref[struct:ultimate-capabilities]{\cmd{ultimate\_capabilities}}.}
{1 when the network target exists.}
{Tests socket-service availability.}
{if (ultimate\_has\_network(\&caps)) connect\_peer();}
{Network code runs only when its target was detected.}

\cmacroref{ultimate\_has\_control}{ultimate\_has\_control(caps)}
{\cmd{caps} points to detected
\hyperref[struct:ultimate-capabilities]{\cmd{ultimate\_capabilities}}.}
{1 when the control target exists.}
{Tests machine-control service availability.}
{if (ultimate\_has\_control(\&caps)) read\_drives();}
{Control operations are attempted only when supported.}

\cmacroref{ultimate\_has\_softiec}{ultimate\_has\_softiec(caps)}
{\cmd{caps} points to detected
\hyperref[struct:ultimate-capabilities]{\cmd{ultimate\_capabilities}}.}
{1 when the SoftwareIEC target identifies itself.}
{Reports target presence, not whether a particular fast load will work;
\cmd{ultimate\_load()} performs the honest command-level test.}
{uint8\_t present = ultimate\_has\_softiec(\&caps);}
{\cmd{present} describes detection only.}

\cmacroref{ultimate\_has\_http}{ultimate\_has\_http(caps)}
{\cmd{caps} points to detected
\hyperref[struct:ultimate-capabilities]{\cmd{ultimate\_capabilities}}.}
{1 when the HTTP target exists.}
{Tests firmware HTTP-client availability.}
{if (ultimate\_has\_http(\&caps)) fetch\_page();}
{HTTP code runs only when its target was detected.}

\section{Palette, Turbo And Vsprite Graphics}

\capiref{ultimate\_palette\_get}{ultimate.h}
{uint8\_t ultimate\_palette\_get(uint8\_t *palette)}
{\cmd{palette} points to \cmd{ULTIMATE\_PALETTE\_BYTES} writable bytes.}
{\cmd{ULTIMATE\_OK}, \cmd{ULTIMATE\_ERR\_NOT\_SUPPORTED}, or
\cmd{ULTIMATE\_ERR\_PROTOCOL} unless a full 48-byte reply arrives.}
{Reads sixteen live RGB triples.}
{uint8\_t palette[ULTIMATE\_PALETTE\_BYTES]; err = ultimate\_palette\_get(palette);}
{\cmd{palette[0..2]} is colour zero's red, green and blue.}

\capiref{ultimate\_palette\_set}{ultimate.h}
{uint8\_t ultimate\_palette\_set(const uint8\_t *palette)}
{\cmd{palette} points to 48 bytes containing sixteen RGB triples.}
{\cmd{ULTIMATE\_OK} or \cmd{ULTIMATE\_ERR\_NOT\_SUPPORTED}.}
{Replaces the live palette in one command without writing flash or a VPL file.}
{err = ultimate\_palette\_set(palette);}
{All sixteen live colours change when the call succeeds.}

\capiref{ultimate\_palette\_set\_color}{ultimate.h}
{uint8\_t ultimate\_palette\_set\_color(uint8\_t index, uint8\_t r, uint8\_t g, uint8\_t b)}
{\cmd{index} is 0--15; RGB components are 0--255.}
{\cmd{ULTIMATE\_OK}, \cmd{ULTIMATE\_ERR\_INVALID\_ARGUMENT} for another index,
or \cmd{ULTIMATE\_ERR\_NOT\_SUPPORTED}.}
{Changes one live palette colour.}
{err = ultimate\_palette\_set\_color(6, 255, 0, 0);}
{Colour index 6 becomes red.}

\capiref{ultimate\_palette\_reset}{ultimate.h}
{uint8\_t ultimate\_palette\_reset(void)}
{None.}
{\cmd{ULTIMATE\_OK} or \cmd{ULTIMATE\_ERR\_NOT\_SUPPORTED}.}
{Restores the machine's built-in live palette.}
{err = ultimate\_palette\_reset();}
{Temporary palette changes are removed.}

\capiref{ultimate\_turbo\_available}{ultimate.h}
{uint8\_t ultimate\_turbo\_available(void)}
{None; \cmd{ultimate\_init()} is not required.}
{1 when the Ultimate~64 turbo registers answer, otherwise 0.}
{Tests the owner's Turbo Control setting through memory-mapped hardware.}
{if (!ultimate\_turbo\_available()) run\_at\_one\_mhz();}
{The fallback runs on a plain C64 or when turbo registers are disabled.}

\capiref{ultimate\_turbo\_get}{ultimate.h}
{uint8\_t ultimate\_turbo\_get(void)}
{None.}
{The current speed index, or \cmd{U64\_TURBO\_UNAVAILABLE} (\$FF).}
{Reads the CPU-speed index without changing the badline setting.}
{uint8\_t old = ultimate\_turbo\_get();}
{\cmd{old} can later restore the entry speed when it is not \$FF.}

\capiref{ultimate\_turbo\_set}{ultimate.h}
{uint8\_t ultimate\_turbo\_set(uint8\_t index)}
{\cmd{index} is 0--15. Constants through \cmd{U64\_SPEED\_4MHZ} are portable;
\cmd{U64\_SPEED\_MAX} means this machine's maximum.}
{\cmd{ULTIMATE\_OK}, \cmd{ULTIMATE\_ERR\_INVALID\_ARGUMENT}, or
\cmd{ULTIMATE\_ERR\_NOT\_SUPPORTED}.}
{Sets CPU speed and leaves badlines unchanged.}
{err = ultimate\_turbo\_set(U64\_SPEED\_1MHZ);\
err = ultimate\_turbo\_set(U64\_SPEED\_2MHZ);\
err = ultimate\_turbo\_set(U64\_SPEED\_3MHZ);\
err = ultimate\_turbo\_set(U64\_SPEED\_4MHZ);\
err = ultimate\_turbo\_set(U64\_SPEED\_MAX);}
{The calls demonstrate every named portable speed and the model's maximum;
unnamed indices 4--14 remain model-dependent.}

\capiref{ultimate\_turbo\_badlines}{ultimate.h}
{uint8\_t ultimate\_turbo\_badlines(uint8\_t on)}
{Zero disables badlines; any nonzero value restores them.}
{\cmd{ULTIMATE\_OK} or \cmd{ULTIMATE\_ERR\_NOT\_SUPPORTED}.}
{Changes VIC badline cycle stealing without changing CPU speed. Check the real
display before shipping with badlines disabled.}
{err = ultimate\_turbo\_badlines(0); err = ultimate\_turbo\_badlines(1);}
{The two calls demonstrate both settings.}

\capiref{ultimate\_vsprite\_draw}{ultimate.h}
{uint8\_t ultimate\_vsprite\_draw(const ultimate\_vsprite *sprite)}
{\cmd{sprite} points to a complete
\hyperref[struct:ultimate-vsprite]{\cmd{ultimate\_vsprite}} descriptor; its
linked type entry defines every field, layout and valid operation.}
{\cmd{ULTIMATE\_OK} or \cmd{ULTIMATE\_ERR\_INVALID\_ARGUMENT}.}
{Composites an OR, masked, bitmap-copy, or colour-aware software vsprite. It
needs no Ultimate or UCI, executes from RAM, self-modifies, and is not
re-entrant.}
{sprite.flags = 0;                 err = ultimate\_vsprite\_draw(\&sprite);\
sprite.flags = VSPRITE\_F\_MASKED; err = ultimate\_vsprite\_draw(\&sprite);\
sprite.flags = VSPRITE\_F\_COPY;   err = ultimate\_vsprite\_draw(\&sprite);\
sprite.flags = VSPRITE\_F\_COLOR;  err = ultimate\_vsprite\_draw(\&sprite);}
{The four calls demonstrate OR, masked, bitmap-copy and colour-aware draws.}

\section{Directories And Files}

\capiref{ultimate\_chdir}{ultimate.h}
{uint8\_t ultimate\_chdir(const char *path)}
{\cmd{path} is a NUL-terminated absolute or relative path.}
{An \cmd{ULTIMATE\_*} result.}
{Changes the current Ultimate DOS directory.}
{err = ultimate\_chdir("/USB1/GAMES");}
{Later filesystem calls use that directory.}

\capiref{ultimate\_getpath}{ultimate.h}
{uint8\_t ultimate\_getpath(char *buf, uint16\_t buflen, uint16\_t *outlen)}
{\cmd{buf} has \cmd{buflen} bytes; \cmd{outlen} may be \cmd{NULL}.}
{An \cmd{ULTIMATE\_*} result; clipping to \cmd{buflen-1} is still success.}
{Returns the current path terminated, and optionally its stored length.}
{err = ultimate\_getpath(path, sizeof path, \&length);}
{\cmd{path} contains the current directory.}

\capiref{ultimate\_opendir}{ultimate.h}
{uint8\_t ultimate\_opendir(void)}
{None.}
{An \cmd{ULTIMATE\_*} result.}
{Begins a block-by-block walk of the current directory. Issue no unrelated
command until it reaches \cmd{ULTIMATE\_END} or is aborted.}
{err = ultimate\_opendir();}
{A successful call prepares the first \cmd{ultimate\_readdir()}.}

\capiref{ultimate\_readdir}{ultimate.h}
{uint8\_t ultimate\_readdir(char *name, uint16\_t namelen, uint8\_t *attrib)}
{\cmd{name} has \cmd{namelen} bytes; \cmd{attrib} may be \cmd{NULL}.}
{\cmd{ULTIMATE\_OK} for one entry, \cmd{ULTIMATE\_END} after the last, or an
error.}
{Consumes one directory reply block, terminates the name, and optionally stores
\cmd{DOS\_ATTR\_*}. Call \cmd{uci\_abort()} when stopping early.}
{err = ultimate\_readdir(name, sizeof name, \&attrib);}
{On \cmd{ULTIMATE\_OK}, \cmd{name} and \cmd{attrib} describe exactly one
entry; call again until \cmd{ULTIMATE\_END}.}

\capiref{ultimate\_open}{ultimate.h}
{uint8\_t ultimate\_open(const char *name, uint8\_t attrib)}
{\cmd{name} is terminated; \cmd{attrib} combines \cmd{DOS\_FA\_READ},
\cmd{DOS\_FA\_WRITE}, \cmd{DOS\_FA\_CREATE\_NEW}, or
\cmd{DOS\_FA\_CREATE\_ALWAYS}.}
{An \cmd{ULTIMATE\_*} result.}
{Opens a firmware-side current file. Create-always replaces; create-new fails
when the name already exists.}
{err = ultimate\_open("INPUT", DOS\_FA\_READ);\
err = ultimate\_open("OUTPUT", DOS\_FA\_WRITE);\
err = ultimate\_open("BOTH", DOS\_FA\_READ | DOS\_FA\_WRITE);\
err = ultimate\_open("NEW", DOS\_FA\_WRITE | DOS\_FA\_CREATE\_NEW);\
err = ultimate\_open("REPLACE", DOS\_FA\_WRITE | DOS\_FA\_CREATE\_ALWAYS);}
{The calls demonstrate read, write, read/write, create-new and
create-or-replace modes.}

\capiref{ultimate\_close}{ultimate.h}
{uint8\_t ultimate\_close(void)}
{None.}
{An \cmd{ULTIMATE\_*} result; closing with nothing open is harmless.}
{Closes the current firmware-side file.}
{err = ultimate\_close();}
{The current file handle is released.}

\capiref{ultimate\_read}{ultimate.h}
{uint8\_t ultimate\_read(uint8\_t *buf, uint16\_t len, uint16\_t *outlen)}
{\cmd{buf} has \cmd{len} bytes; \cmd{outlen} may be \cmd{NULL}; a file is open.}
{An \cmd{ULTIMATE\_*} result. A short successful read means end-of-file.}
{Reads up to \cmd{len} bytes, walking the firmware's reply blocks internally.}
{err = ultimate\_read(buf, sizeof buf, \&got);}
{\cmd{got} is the actual byte count when supplied.}

\capiref{ultimate\_write}{ultimate.h}
{uint8\_t ultimate\_write(const uint8\_t *buf, uint16\_t len)}
{\cmd{buf} points to \cmd{len} bytes; a writable file is open.}
{An \cmd{ULTIMATE\_*} result.}
{Writes the complete caller span to the current file.}
{err = ultimate\_write(message, sizeof message);}
{The bytes are appended or written at the current file position.}

\capiref{ultimate\_seek}{ultimate.h}
{uint8\_t ultimate\_seek(uint32\_t pos)}
{\cmd{pos} is an absolute byte position in the current open file.}
{An \cmd{ULTIMATE\_*} result.}
{Moves the firmware-side file position.}
{err = ultimate\_seek(1024UL);}
{The next read or write begins at byte 1024.}

\capiref{ultimate\_delete}{ultimate.h}
{uint8\_t ultimate\_delete(const char *name)}
{\cmd{name} is a terminated current-directory name or path.}
{An \cmd{ULTIMATE\_*} result.}
{Deletes one file.}
{err = ultimate\_delete("RESULT.TXT");}
{RESULT.TXT is removed on success.}

\capiref{ultimate\_stat}{ultimate.h}
{uint8\_t ultimate\_stat(const char *name, ultimate\_fileinfo *info)}
{\cmd{name} is terminated; \cmd{info} points to writable
\hyperref[struct:ultimate-fileinfo]{\cmd{ultimate\_fileinfo}} storage.}
{\cmd{ULTIMATE\_OK}, \cmd{ULTIMATE\_ERR\_PROTOCOL} for a short reply, or
another error.}
{Reads size, packed FAT date/time, extension, attributes and terminated name.}
{ultimate\_fileinfo info; err = ultimate\_stat("DEMO.PRG", \&info);}
{\cmd{info.size} is valid only after success.}

\capiref{ultimate\_fstat}{ultimate.h}
{uint8\_t ultimate\_fstat(ultimate\_fileinfo *info)}
{\cmd{info} points to writable
\hyperref[struct:ultimate-fileinfo]{\cmd{ultimate\_fileinfo}} storage; a file
is open.}
{\cmd{ULTIMATE\_OK}, \cmd{ULTIMATE\_ERR\_PROTOCOL}, or another error.}
{Reads metadata for the current open file.}
{ultimate\_fileinfo info; err = ultimate\_fstat(\&info);}
{The open file's metadata fills \cmd{info}.}

\capiref{ultimate\_rename}{ultimate.h}
{uint8\_t ultimate\_rename(const char *from, const char *to)}
{Both pointers name terminated entries in the current directory.}
{An \cmd{ULTIMATE\_*} result.}
{Renames one entry without copying its contents through the C64.}
{err = ultimate\_rename("OLD.TXT", "NEW.TXT");}
{OLD.TXT becomes NEW.TXT.}

\capiref{ultimate\_copy}{ultimate.h}
{uint8\_t ultimate\_copy(const char *name, const char *destination\_path)}
{\cmd{name} is a current-directory source; \cmd{destination\_path} names the
target directory.}
{An \cmd{ULTIMATE\_*} result.}
{Copies inside the Ultimate and preserves the source filename.}
{err = ultimate\_copy("NEW.TXT", "/USB1/BACKUP");}
{The destination is /USB1/BACKUP/NEW.TXT.}

\capiref{ultimate\_mkdir}{ultimate.h}
{uint8\_t ultimate\_mkdir(const char *name)}
{\cmd{name} is a terminated new directory name.}
{An \cmd{ULTIMATE\_*} result.}
{Creates one directory in the current directory.}
{err = ultimate\_mkdir("SAVES");}
{SAVES is created.}

\capiref{ultimate\_home}{ultimate.h}
{uint8\_t ultimate\_home(void)}
{None.}
{\cmd{ULTIMATE\_OK} or \cmd{ULTIMATE\_ERR\_NOT\_SUPPORTED} when no home is
configured.}
{Changes to the Ultimate's configured home directory.}
{err = ultimate\_home();}
{Later filesystem calls use the configured home.}

\section{Disk Images And Clock}

\capiref{ultimate\_mount}{ultimate.h}
{uint8\_t ultimate\_mount(uint8\_t device, const char *image)}
{\cmd{device} is an IEC number or \cmd{ULTIMATE\_DRIVE\_LAST}; \cmd{image} is a
terminated D64, D71, D81, G64 or G71 path.}
{An \cmd{ULTIMATE\_*} result; a disabled drive normally reports
\cmd{ULTIMATE\_ERR\_DEVICE}.}
{Mounts an image in the selected emulated drive.}
{err = ultimate\_mount(8, "GAMES/PITFALL.D64");}
{PITFALL.D64 is mounted in the drive answering as IEC device 8.}

\capiref{ultimate\_unmount}{ultimate.h}
{uint8\_t ultimate\_unmount(uint8\_t device)}
{\cmd{device} is an IEC number or \cmd{ULTIMATE\_DRIVE\_LAST}.}
{An \cmd{ULTIMATE\_*} result.}
{Removes the current image from a drive.}
{err = ultimate\_unmount(8);}
{The drive answering as device 8 becomes empty.}

\capiref{ultimate\_swap}{ultimate.h}
{uint8\_t ultimate\_swap(uint8\_t device)}
{\cmd{device} is an IEC number or \cmd{ULTIMATE\_DRIVE\_LAST}.}
{An \cmd{ULTIMATE\_*} result.}
{Advances a multi-image set the same way as the Ultimate menu.}
{err = ultimate\_swap(8);}
{Device 8 selects its next image.}

\capiref{ultimate\_get\_time}{ultimate.h}
{uint8\_t ultimate\_get\_time(uint8\_t format, char *buf, uint16\_t buflen, uint16\_t *outlen)}
{\cmd{format} is \cmd{ULTIMATE\_TIME\_PLAIN} or
\cmd{ULTIMATE\_TIME\_WEEKDAY}; \cmd{buf} has \cmd{buflen} bytes;
\cmd{outlen} may be \cmd{NULL}.}
{An \cmd{ULTIMATE\_*} result.}
{Reads the battery-backed clock as terminated firmware-formatted text.}
{err = ultimate\_get\_time(ULTIMATE\_TIME\_PLAIN, now, sizeof now, \&length);\
err = ultimate\_get\_time(ULTIMATE\_TIME\_WEEKDAY, now, sizeof now, \&length);}
{Plain output resembles \cmd{2026/08/22 14:30:00}; weekday output prefixes a
day name.}

\capiref{ultimate\_set\_time}{ultimate.h}
{uint8\_t ultimate\_set\_time(uint8\_t year\_1900, uint8\_t month, uint8\_t day, uint8\_t hour, uint8\_t minute, uint8\_t second)}
{The year is calendar year minus 1900; remaining arguments are month, day,
hour, minute and second.}
{An \cmd{ULTIMATE\_*} result.}
{Sets the battery-backed clock.}
{err = ultimate\_set\_time(126, 8, 22, 14, 30, 0);}
{The clock becomes 2026-08-22 14:30:00.}

\section{Loading And Saving Memory}

\capiref{ultimate\_load}{ultimate.h}
{uint8\_t ultimate\_load(const char *name, uint16\_t addr)}
{\cmd{name} is a non-null terminated PRG path; zero \cmd{addr} uses its header,
and a nonzero value overrides it.}
{An \cmd{ULTIMATE\_*} result; after success
\cmd{ultimate\_last\_end()} supplies the end address.}
{Consumes the two-byte PRG header and loads its body. The SDK first tries
SoftwareIEC and transparently falls back to DOS reads.}
{err = ultimate\_load("GAME.PRG", 0); err = ultimate\_load("FONT.PRG", 0xc000);}
{The calls demonstrate header-selected and caller-selected addresses.}

\capiref{ultimate\_bload}{ultimate.h}
{uint8\_t ultimate\_bload(const char *name, uint16\_t addr, uint16\_t max)}
{\cmd{name} is terminated; \cmd{addr} and \cmd{max} must be nonzero.}
{An \cmd{ULTIMATE\_*} result; \cmd{ultimate\_last\_end()} supplies the end
address after success.}
{Loads at most \cmd{max} raw bytes without consuming a header.}
{err = ultimate\_bload("FONT.BIN", 0xc000, 2048);}
{At most 2048 bytes are written from \$C000.}

\capiref{ultimate\_save}{ultimate.h}
{uint8\_t ultimate\_save(const char *name, uint16\_t start, uint16\_t len)}
{\cmd{name} is terminated; \cmd{start} is a C64 address; \cmd{len} is nonzero.}
{An \cmd{ULTIMATE\_*} result.}
{Creates or replaces a file with a C64 memory range.}
{err = ultimate\_save("SCREEN.BIN", 0x0400, 1000);}
{SCREEN.BIN receives the 1000-byte text screen.}

\capiref{ultimate\_last\_end}{ultimate.h}
{uint16\_t ultimate\_last\_end(void)}
{None.}
{The address immediately following the last successful
\cmd{ultimate\_load()} or \cmd{ultimate\_bload()}.}
{Retrieves cached load state without touching hardware.}
{uint16\_t end = ultimate\_last\_end();}
{\cmd{end-1} is the final address written by the preceding successful load.}

\section{RAM Expansion}

\capiref{ultimate\_reu\_available}{ultimate.h}
{uint8\_t ultimate\_reu\_available(void)}
{None; UCI initialization is not required.}
{1 when the REU registers answer, otherwise 0.}
{Tests whether RAM expansion is enabled.}
{if (!ultimate\_reu\_available()) return ULTIMATE\_ERR\_NOT\_SUPPORTED;}
{The program refuses an REU-dependent path when no expansion is mapped.}

\capiref{ultimate\_reu\_size}{ultimate.h}
{uint16\_t ultimate\_reu\_size(void)}
{None.}
{Size in 64K banks, or zero when unavailable.}
{Measures the expansion non-destructively by testing alias boundaries.}
{uint16\_t banks = ultimate\_reu\_size();}
{2 means 128\,KB, 16 means 1\,MB, 128 means 8\,MB and 256 means 16\,MB.}

\capiref{ultimate\_reu\_stash}{ultimate.h}
{uint8\_t ultimate\_reu\_stash(uint16\_t addr, uint32\_t reuaddr, uint16\_t len)}
{\cmd{addr} is a C64 address; \cmd{reuaddr} is within 24 bits; zero
\cmd{len} means 65536 bytes.}
{\cmd{ULTIMATE\_OK}, \cmd{ULTIMATE\_ERR\_NOT\_SUPPORTED}, or
\cmd{ULTIMATE\_ERR\_INVALID\_ARGUMENT}.}
{DMA-copies C64 memory into the expansion and returns after completion.}
{err = ultimate\_reu\_stash(0x0400, 0, 1000);}
{The text screen is copied to REU address zero.}

\capiref{ultimate\_reu\_fetch}{ultimate.h}
{uint8\_t ultimate\_reu\_fetch(uint16\_t addr, uint32\_t reuaddr, uint16\_t len)}
{\cmd{addr} is a C64 address; \cmd{reuaddr} is within 24 bits; zero
\cmd{len} means 65536 bytes.}
{\cmd{ULTIMATE\_OK}, \cmd{ULTIMATE\_ERR\_NOT\_SUPPORTED}, or
\cmd{ULTIMATE\_ERR\_INVALID\_ARGUMENT}.}
{DMA-copies expansion data into C64 memory.}
{err = ultimate\_reu\_fetch(0x0400, 0, 1000);}
{The first 1000 REU bytes replace the text screen.}

\capiref{ultimate\_reu\_load}{ultimate.h}
{uint8\_t ultimate\_reu\_load(uint32\_t reuaddr, uint32\_t len)}
{A readable file is open; \cmd{reuaddr} and \cmd{len} describe a range within
the 16\,MB REU window.}
{An \cmd{ULTIMATE\_*} result.}
{Transfers the current open file directly into the expansion. It may take a
long time and restores the caller's UCI timeout.}
{err = ultimate\_open("ASSETS.BIN", DOS\_FA\_READ); err = ultimate\_reu\_load(0, 65536UL);}
{The first 65536 file bytes move to REU address zero without passing through C64 RAM.}

\capiref{ultimate\_reu\_save}{ultimate.h}
{uint8\_t ultimate\_reu\_save(uint32\_t reuaddr, uint32\_t len)}
{A writable file is open; \cmd{reuaddr} and \cmd{len} describe a range within
the 16\,MB REU window.}
{An \cmd{ULTIMATE\_*} result.}
{Transfers expansion data directly into the current open file and restores the
caller's timeout.}
{err = ultimate\_open("CAPTURE.BIN", DOS\_FA\_WRITE | DOS\_FA\_CREATE\_ALWAYS); err = ultimate\_reu\_save(0, 65536UL);}
{65536 REU bytes are written without passing through C64 RAM.}

\begin{tip}
Keep large level data, decompression sources, sampled audio and screen backups
in the REU. Use stash/fetch for small live-memory swaps and load/save when the
data should move directly between storage and expansion memory.
\end{tip}

\section{Ultimate Audio}

\capiref{ultimate\_audio\_init}{ultimate.h}
{uint8\_t ultimate\_audio\_init(void)}
{Ultimate Audio must be mapped at \$DF20--\$DFFF.}
{\cmd{ULTIMATE\_OK} or \cmd{ULTIMATE\_ERR\_NOT\_SUPPORTED}.}
{Probes the PCM engine with a silent byte. Call once before other audio calls.}
{err = ultimate\_audio\_init();}
{Audio services are usable only after success.}

\capiref{ultimate\_audio\_available}{ultimate.h}
{uint8\_t ultimate\_audio\_available(void)}
{Call \cmd{ultimate\_audio\_init()} first.}
{1 after successful audio initialization, otherwise 0.}
{Reads cached audio state without touching hardware.}
{if (!ultimate\_audio\_available()) return 1;}
{The program exits when the PCM engine is unavailable.}

\capiref{ultimate\_audio\_version}{ultimate.h}
{uint8\_t ultimate\_audio\_version(void)}
{Audio must be available.}
{The raw module-version byte.}
{Reads the mapped engine's version register.}
{uint8\_t version = ultimate\_audio\_version();}
{\cmd{version} is meaningful only after audio initialization succeeds.}

\capiref{ultimate\_audio\_configure}{ultimate.h}
{uint8\_t ultimate\_audio\_configure(const ultimate\_audio\_voice *voice)}
{\cmd{voice} points to a fully initialized
\hyperref[struct:ultimate-audio-voice]{\cmd{ultimate\_audio\_voice}}; its linked
type entry defines all fields and constraints.}
{\cmd{ULTIMATE\_OK}, \cmd{ULTIMATE\_ERR\_INVALID\_ARGUMENT}, or
\cmd{ULTIMATE\_ERR\_NOT\_SUPPORTED}.}
{Stops, validates and programs one voice, leaving its gate closed. Flags select
8/16-bit, mono/interleaved, repeat and IRQ operation.}
{voice.flags = 0;                  err = ultimate\_audio\_configure(\&voice);\
voice.flags = UA\_CTRL\_16BIT;    err = ultimate\_audio\_configure(\&voice);\
voice.flags = UA\_CTRL\_INTERLEAVE; err = ultimate\_audio\_configure(\&voice);\
voice.flags = UA\_CTRL\_REPEAT;   err = ultimate\_audio\_configure(\&voice);\
voice.flags = UA\_CTRL\_IRQ;      err = ultimate\_audio\_configure(\&voice);}
{The calls isolate 8-bit mono, 16-bit, stereo interleave, repeat and completion
IRQ settings; valid flags may be combined.}

\capiref{ultimate\_audio\_start}{ultimate.h}
{uint8\_t ultimate\_audio\_start(uint8\_t channel, uint8\_t flags)}
{\cmd{channel} is 0 through \cmd{UA\_CHANNELS-1}; \cmd{flags} repeats only the
configured \cmd{UA\_CTRL\_FLAGS} format/playback bits. Do not include
\cmd{UA\_CTRL\_GATE}; the call adds it.}
{\cmd{ULTIMATE\_OK}, \cmd{ULTIMATE\_ERR\_INVALID\_ARGUMENT}, or
\cmd{ULTIMATE\_ERR\_NOT\_SUPPORTED}.}
{Opens a configured channel's gate immediately.}
{err = ultimate\_audio\_start(0, UA\_CTRL\_16BIT | UA\_CTRL\_REPEAT);}
{Channel zero begins playing.}

\capiref{ultimate\_audio\_stop}{ultimate.h}
{uint8\_t ultimate\_audio\_stop(uint8\_t channel)}
{\cmd{channel} is 0 through \cmd{UA\_CHANNELS-1}.}
{\cmd{ULTIMATE\_OK}, \cmd{ULTIMATE\_ERR\_INVALID\_ARGUMENT}, or
\cmd{ULTIMATE\_ERR\_NOT\_SUPPORTED}.}
{Closes one channel's gate immediately.}
{err = ultimate\_audio\_stop(0);}
{Channel zero stops.}

\capiref{ultimate\_audio\_irq\_status}{ultimate.h}
{uint8\_t ultimate\_audio\_irq\_status(void)}
{Audio must be available.}
{A bit mask with one end-of-sample bit per channel.}
{Reads latched PCM completion state.}
{uint8\_t ended = ultimate\_audio\_irq\_status();}
{\cmd{ended \& (1u << channel)} is nonzero when that channel completed.}

\capiref{ultimate\_audio\_irq\_clear}{ultimate.h}
{uint8\_t ultimate\_audio\_irq\_clear(uint8\_t channel)}
{\cmd{channel} is 0 through \cmd{UA\_CHANNELS-1}.}
{\cmd{ULTIMATE\_OK}, \cmd{ULTIMATE\_ERR\_INVALID\_ARGUMENT}, or
\cmd{ULTIMATE\_ERR\_NOT\_SUPPORTED}.}
{Clears one channel's latched completion bit and its asserted IRQ.}
{err = ultimate\_audio\_irq\_clear(channel);}
{That channel's completion bit is clear.}

\capiref{ultimate\_audio\_load\_wav}{ultimate.h}
{uint8\_t ultimate\_audio\_load\_wav(const char *name, uint32\_t reuaddr, ultimate\_audio\_voice *voice)}
{\cmd{name} is a terminated PCM WAV path; \cmd{reuaddr} begins its REU data;
\cmd{voice} points to
\hyperref[struct:ultimate-audio-voice]{\cmd{ultimate\_audio\_voice}} and
preserves caller channel, volume, pan and repeat points.}
{DOS errors, \cmd{ULTIMATE\_ERR\_PROTOCOL} for broken RIFF,
\cmd{ULTIMATE\_ERR\_NOT\_SUPPORTED} for unsupported format/rate/range, or an
REU error. The voice is invalid on failure.}
{Loads 8/16-bit mono/stereo PCM, fixes unsigned 8-bit samples in place, closes
the file on every path, and fills address, length, rate and format flags.}
{err = ultimate\_audio\_load\_wav("/USB1/BOING.WAV", 0x4000UL, \&voice);}
{On success, configure and start \cmd{voice}; for stereo derive the right
channel as described in Chapter~3.}

\section{Machine And Drive Control}

\capiref{ultimate\_reboot}{ultimate.h}
{uint8\_t ultimate\_reboot(void)}
{None.}
{Does not return.}
{Resets the C64; the calling program stops existing during the command.}
{ultimate\_reboot(); /* no following statement runs */}
{The machine restarts.}

\capiref{ultimate\_freeze}{ultimate.h}
{uint8\_t ultimate\_freeze(void)}
{None.}
{An \cmd{ULTIMATE\_*} result after the user leaves the menu.}
{Invokes the Ultimate freeze menu and waits for human release.}
{err = ultimate\_freeze();}
{Execution resumes after the menu closes.}

\capiref{ultimate\_drive\_enable}{ultimate.h}
{uint8\_t ultimate\_drive\_enable(uint8\_t drive, uint8\_t on)}
{\cmd{drive} is \cmd{ULTIMATE\_DRIVE\_A} or \cmd{ULTIMATE\_DRIVE\_B}; zero
\cmd{on} disables, nonzero enables.}
{An \cmd{ULTIMATE\_*} result.}
{Switches one physical emulated drive off or on.}
{err = ultimate\_drive\_enable(ULTIMATE\_DRIVE\_A, 0); err = ultimate\_drive\_enable(ULTIMATE\_DRIVE\_A, 1);}
{The calls demonstrate both states.}

\capiref{ultimate\_drive\_power}{ultimate.h}
{uint8\_t ultimate\_drive\_power(uint8\_t drive, uint8\_t *on)}
{\cmd{drive} is A or B; \cmd{on} is a required writable pointer.}
{An \cmd{ULTIMATE\_*} result; \cmd{*on} is 1 while running.}
{Reads one physical drive's power state.}
{err = ultimate\_drive\_power(ULTIMATE\_DRIVE\_A, \&running);}
{\cmd{running} contains zero or one after success.}

\capiref{ultimate\_drive\_info}{ultimate.h}
{uint8\_t ultimate\_drive\_info(ultimate\_drives *drives)}
{\cmd{drives} points to writable
\hyperref[struct:ultimate-drives]{\cmd{ultimate\_drives}} storage.}
{An \cmd{ULTIMATE\_*} result; count never exceeds records actually received.}
{Reports type, IEC device number and power for disk drives and occupied IEC
slots. SoftwareIEC and printer records are bus slots, not physical drives.}
{ultimate\_drives drives; err = ultimate\_drive\_info(\&drives);}
{\cmd{drives.drive[i].device} is the number accepted by mount operations.}

\capiref{ultimate\_ramdisk\_info}{ultimate.h}
{uint8\_t ultimate\_ramdisk\_info(uint8\_t *info)}
{\cmd{info} has \cmd{CTRL\_RAMDISK\_BYTES} writable bytes.}
{An \cmd{ULTIMATE\_*} result.}
{Reads GEOS MegaPatch RAM-disk records: type followed by size in 64K units.}
{uint8\_t info[CTRL\_RAMDISK\_BYTES]; err = ultimate\_ramdisk\_info(info);}
{A record type of zero marks an empty slot.}

\section{Network Sockets}

\capiref{ultimate\_net\_ifcount}{ultimate.h}
{uint8\_t ultimate\_net\_ifcount(uint8\_t *count)}
{\cmd{count} is a required writable pointer.}
{An \cmd{ULTIMATE\_*} result; \cmd{*count} is valid after success.}
{Reports the number of network interfaces.}
{err = ultimate\_net\_ifcount(\&count);}
{Indices zero through \cmd{count-1} may be queried.}

\capiref{ultimate\_net\_macaddr}{ultimate.h}
{uint8\_t ultimate\_net\_macaddr(uint8\_t iface, uint8\_t *mac)}
{\cmd{iface} is an interface index; \cmd{mac} has
\cmd{UCI\_NET\_MACADDR\_BYTES} writable bytes.}
{An \cmd{ULTIMATE\_*} result.}
{Reads the six-byte hardware address.}
{uint8\_t mac[UCI\_NET\_MACADDR\_BYTES]; err = ultimate\_net\_macaddr(0, mac);}
{\cmd{mac[0..5]} contains interface zero's address.}

\capiref{ultimate\_net\_ipconfig}{ultimate.h}
{uint8\_t ultimate\_net\_ipconfig(uint8\_t iface, uint8\_t *ipconfig)}
{\cmd{iface} is an interface index; \cmd{ipconfig} has 12 writable bytes.}
{An \cmd{ULTIMATE\_*} result.}
{Reads IPv4 address, netmask and gateway as three four-byte network-order
addresses.}
{uint8\_t ip[UCI\_NET\_IPCONFIG\_BYTES]; err = ultimate\_net\_ipconfig(0, ip);}
{\cmd{ip[0..3]}, \cmd{ip[4..7]} and \cmd{ip[8..11]} hold the three addresses.}

\capiref{ultimate\_net\_setip}{ultimate.h}
{uint8\_t ultimate\_net\_setip(uint8\_t iface, const uint8\_t *ipconfig)}
{\cmd{iface} is an interface index; \cmd{ipconfig} supplies the same 12-byte
address/netmask/gateway layout returned above.}
{An \cmd{ULTIMATE\_*} result.}
{Changes only the running network configuration; it does not write stored
settings or detect address conflicts.}
{err = ultimate\_net\_setip(0, ip);}
{Interface zero immediately uses the supplied configuration.}

\capiref{ultimate\_net\_connect}{ultimate.h}
{uint8\_t ultimate\_net\_connect(const char *host, uint16\_t port, uint8\_t *handle)}
{\cmd{host} is terminated; \cmd{port} is a TCP port; \cmd{handle} is required.}
{An \cmd{ULTIMATE\_*} result; \cmd{*handle} receives the socket on success.}
{Opens a TCP socket. DNS/connect may take about 30 seconds, so this call waits
without the normal timeout limit and restores the caller's budget.}
{err = ultimate\_net\_connect("192.0.2.1", 80, \&handle);}
{\cmd{handle} names the connected TCP socket.}

\capiref{ultimate\_net\_udp}{ultimate.h}
{uint8\_t ultimate\_net\_udp(const char *host, uint16\_t port, uint8\_t *handle)}
{\cmd{host} is terminated; \cmd{port} is a UDP port; \cmd{handle} is required.}
{An \cmd{ULTIMATE\_*} result; \cmd{*handle} receives the socket on success.}
{Opens a connected UDP socket with the same unbounded-connect timeout handling
as the TCP call.}
{err = ultimate\_net\_udp("192.0.2.1", 6502, \&handle);}
{\cmd{handle} names the UDP socket.}

\capiref{ultimate\_net\_close}{ultimate.h}
{uint8\_t ultimate\_net\_close(uint8\_t handle)}
{\cmd{handle} names an open socket.}
{An \cmd{ULTIMATE\_*} result.}
{Closes a socket explicitly. Do not close after \cmd{ultimate\_net\_read()}
returns \cmd{ULTIMATE\_END}; the firmware already removed it.}
{err = ultimate\_net\_close(handle);}
{The socket handle is released.}

\capiref{ultimate\_net\_read}{ultimate.h}
{uint8\_t ultimate\_net\_read(uint8\_t handle, uint8\_t *buf, uint16\_t bufsize, uint16\_t *got)}
{\cmd{handle} is open; \cmd{buf} has \cmd{bufsize} bytes; \cmd{got} is required.
The SDK refuses unsafe requests above \cmd{UCI\_NET\_READ\_MAX}.}
{\cmd{ULTIMATE\_OK} with zero or more bytes, \cmd{ULTIMATE\_END} once when the
peer closes, or an error.}
{Polls without waiting. At most \cmd{bufsize-UCI\_NET\_READ\_PREFIX} data bytes
fit because the firmware prefixes its count.}
{err = ultimate\_net\_read(handle, buf, sizeof buf, \&got);}
{Use data only when \cmd{err == ULTIMATE\_OK \&\& got != 0}; stop on
\cmd{ULTIMATE\_END}.}

\capiref{ultimate\_net\_write}{ultimate.h}
{uint8\_t ultimate\_net\_write(uint8\_t handle, const uint8\_t *buf, uint16\_t len, uint16\_t *sent)}
{\cmd{handle} is open; \cmd{buf} supplies \cmd{len} bytes; \cmd{sent} is
required.}
{An \cmd{ULTIMATE\_*} result; \cmd{*sent} is the accepted count.}
{Writes one span and reports rather than assumes how much the firmware took.}
{err = ultimate\_net\_write(handle, buf, len, \&sent);}
{After success, compare \cmd{sent} with \cmd{len}.}

\begin{tip}
Use a small state machine around nonblocking reads: connect once, poll during
the program's normal update loop, consume available bytes, and treat
\cmd{ULTIMATE\_END} as the terminal state.
\end{tip}

\section{HTTP Requests}

\capiref{ultimate\_http\_get}{ultimate.h}
{uint8\_t ultimate\_http\_get(const char *url, uint8\_t *buf, uint16\_t bufsize, uint16\_t *got)}
{\cmd{url} is terminated ASCII; \cmd{buf} has \cmd{bufsize} bytes;
\cmd{got} is required.}
{\cmd{ULTIMATE\_OK} for HTTP below 400, \cmd{ULTIMATE\_ERR\_DEVICE} for HTTP
errors, \cmd{ULTIMATE\_ERR\_TRUNCATED} when the body is larger, or another
transport error.}
{Creates, executes and frees a GET request. It waits through DNS/connect and
restores the timeout. The HTTP status remains in
\cmd{ultimate\_device\_code()}.}
{err = ultimate\_http\_get(url, buf, sizeof buf, \&got);}
{\cmd{buf[0..got-1]} contains the retained body, including an HTTP error body.}

\capiref{ultimate\_http\_open}{ultimate.h}
{uint8\_t ultimate\_http\_open(uint8\_t verb, const char *url, uint8\_t *handle)}
{\cmd{verb} is any \cmd{HTTP\_VERB\_*} from GET through TRACE; \cmd{url} is
terminated ASCII; \cmd{handle} is required.}
{An \cmd{ULTIMATE\_*} result; \cmd{*handle} receives a request slot.}
{Creates a request for later headers and exchange.}
{err = ultimate\_http\_open(HTTP\_VERB\_GET, url, \&request);\
err = ultimate\_http\_open(HTTP\_VERB\_PUT, url, \&request);\
err = ultimate\_http\_open(HTTP\_VERB\_POST, url, \&request);\
err = ultimate\_http\_open(HTTP\_VERB\_PATCH, url, \&request);\
err = ultimate\_http\_open(HTTP\_VERB\_DELETE, url, \&request);\
err = ultimate\_http\_open(HTTP\_VERB\_HEAD, url, \&request);\
err = ultimate\_http\_open(HTTP\_VERB\_OPTIONS, url, \&request);\
err = ultimate\_http\_open(HTTP\_VERB\_CONNECT, url, \&request);\
err = ultimate\_http\_open(HTTP\_VERB\_TRACE, url, \&request);}
{Each call allocates a separate request using the named HTTP method; close each
successful handle.}

\capiref{ultimate\_http\_header}{ultimate.h}
{uint8\_t ultimate\_http\_header(uint8\_t handle, const char *line)}
{\cmd{handle} is an open request; \cmd{line} is terminated
\cmd{key: value} ASCII.}
{An \cmd{ULTIMATE\_*} result.}
{Adds one HTTP header. Call repeatedly for multiple fields.}
{static const char header[] = \{\
0x61,0x63,0x63,0x65,0x70,0x74,0x3a,0x20,\
0x61,0x70,0x70,0x6c,0x69,0x63,0x61,0x74,0x69,0x6f,0x6e,0x2f,\
0x6a,0x73,0x6f,0x6e,0\};\
err = ultimate\_http\_header(request, header);}
{The header is attached to the request slot.}

\capiref{ultimate\_http\_exchange}{ultimate.h}
{uint8\_t ultimate\_http\_exchange(uint8\_t handle, uint8\_t body, uint8\_t *buf, uint16\_t bufsize, uint16\_t *got)}
{\cmd{handle} is an open request; \cmd{body} is a body handle or
\cmd{HTTP\_BODY\_NONE}; \cmd{buf} has \cmd{bufsize} bytes; \cmd{got} is
required.}
{The same HTTP/result/truncation rules as \cmd{ultimate\_http\_get()}.}
{Sends the prepared request. It waits through DNS/connect and restores the
timeout, but does not free the header or body slots.}
{err = ultimate\_http\_exchange(request, HTTP\_BODY\_NONE, buf, sizeof buf, \&got);}
{A request without a body completes and retains up to \cmd{sizeof buf} bytes.}

\capiref{ultimate\_http\_close}{ultimate.h}
{uint8\_t ultimate\_http\_close(uint8\_t handle)}
{\cmd{handle} names an open request/header slot.}
{An \cmd{ULTIMATE\_*} result.}
{Returns one request slot to the firmware.}
{err = ultimate\_http\_close(request);}
{The request slot is reusable.}

\capiref{ultimate\_http\_free\_all}{ultimate.h}
{uint8\_t ultimate\_http\_free\_all(void)}
{None.}
{An \cmd{ULTIMATE\_*} result.}
{Frees every HTTP header and body slot held by the machine. Use for startup
recovery; it may also invalidate slots owned by another C64 program.}
{err = ultimate\_http\_free\_all();}
{All sixteen firmware slots become available.}

\section{HTTP Request Bodies}

\capiref{ultimate\_http\_body}{ultimate.h}
{uint8\_t ultimate\_http\_body(uint8\_t format, uint8\_t *handle)}
{\cmd{format} is \cmd{HTTP\_BODY\_BINARY}, \cmd{HTTP\_BODY\_JSON\_OBJECT},
\cmd{HTTP\_BODY\_JSON\_ARRAY}, or \cmd{HTTP\_BODY\_URL\_ENCODED};
\cmd{handle} is required.}
{An \cmd{ULTIMATE\_*} result; \cmd{*handle} receives the new slot.}
{Creates an empty firmware-side body in any supported format. Binary bodies
accept only the binary append call; the other formats accept typed items.}
{err = ultimate\_http\_body(HTTP\_BODY\_BINARY, \&body);\
err = ultimate\_http\_body(HTTP\_BODY\_JSON\_OBJECT, \&body);\
err = ultimate\_http\_body(HTTP\_BODY\_JSON\_ARRAY, \&body);\
err = ultimate\_http\_body(HTTP\_BODY\_URL\_ENCODED, \&body);}
{The calls create binary, JSON-object, JSON-array and URL-encoded bodies;
free every successful handle.}

\capiref{ultimate\_http\_body\_free}{ultimate.h}
{uint8\_t ultimate\_http\_body\_free(uint8\_t handle)}
{\cmd{handle} names a body slot.}
{An \cmd{ULTIMATE\_*} result.}
{Returns one body slot to the firmware.}
{err = ultimate\_http\_body\_free(body);}
{The body slot is reusable.}

\capiref{ultimate\_http\_body\_clear}{ultimate.h}
{uint8\_t ultimate\_http\_body\_clear(uint8\_t handle)}
{\cmd{handle} names a body slot.}
{An \cmd{ULTIMATE\_*} result.}
{Empties a body while retaining its slot and format.}
{err = ultimate\_http\_body\_clear(body);}
{The body can be rebuilt without allocating another slot.}

\capiref{ultimate\_http\_body\_up}{ultimate.h}
{uint8\_t ultimate\_http\_body\_up(uint8\_t handle)}
{\cmd{handle} names a JSON/form body currently inside a child container.}
{An \cmd{ULTIMATE\_*} result.}
{Leaves the current object or array and returns to its parent.}
{err = ultimate\_http\_body\_up(body);}
{Subsequent items are added beside the completed child container.}

\capiref{ultimate\_http\_body\_int}{ultimate.h}
{uint8\_t ultimate\_http\_body\_int(uint8\_t handle, const char *key, int32\_t value)}
{\cmd{handle} names a nonbinary body; \cmd{key} is non-NULL and at most
\cmd{ULTIMATE\_HTTP\_KEY\_MAX} bytes. Inside arrays pass an empty string; the
pointed-to key is ignored.}
{An \cmd{ULTIMATE\_*} result, including
\cmd{ULTIMATE\_ERR\_INVALID\_ARGUMENT} for an overlong key.}
{Adds a signed 32-bit integer.}
{static const char key[] = \{0x73,0x63,0x6f,0x72,0x65,0\};\
err = ultimate\_http\_body\_int(body, key, -6502L);}
{An object gets \cmd{"score":-6502}; for an array, pass an empty key string and
only the value is added.}

\capiref{ultimate\_http\_body\_bool}{ultimate.h}
{uint8\_t ultimate\_http\_body\_bool(uint8\_t handle, const char *key, uint8\_t value)}
{\cmd{handle} names a nonbinary body; \cmd{key} follows the same rules; zero is
false and nonzero is true.}
{An \cmd{ULTIMATE\_*} result.}
{Adds a Boolean value.}
{static const char key[] = \{0x72,0x65,0x61,0x64,0x79,0\};\
err = ultimate\_http\_body\_bool(body, key, 1);\
err = ultimate\_http\_body\_bool(body, key, 0);}
{The calls add true and false values respectively.}

\capiref{ultimate\_http\_body\_string}{ultimate.h}
{uint8\_t ultimate\_http\_body\_string(uint8\_t handle, const char *key, const char *value)}
{\cmd{handle} names a nonbinary body; \cmd{key} follows the same rules;
\cmd{value} is terminated and at most 255 bytes.}
{An \cmd{ULTIMATE\_*} result.}
{Adds one string value.}
{static const char key[] = \{0x6e,0x61,0x6d,0x65,0\};\
static const char value[] = \{0x62,0x6f,0x69,0x6e,0x67,0\};\
err = ultimate\_http\_body\_string(body, key, value);}
{The object gets \cmd{"name":"boing"}.}

\capiref{ultimate\_http\_body\_object}{ultimate.h}
{uint8\_t ultimate\_http\_body\_object(uint8\_t handle, const char *key)}
{\cmd{handle} names a nonbinary body; \cmd{key} follows the same rules.}
{An \cmd{ULTIMATE\_*} result.}
{Adds an object and enters it. Call \cmd{ultimate\_http\_body\_up()} to leave.}
{static const char key[] = \{0x70,0x6c,0x61,0x79,0x65,0x72,0\};\
err = ultimate\_http\_body\_object(body, key);}
{Following typed values are children of \cmd{player}.}

\capiref{ultimate\_http\_body\_array}{ultimate.h}
{uint8\_t ultimate\_http\_body\_array(uint8\_t handle, const char *key)}
{\cmd{handle} names a nonbinary body; \cmd{key} follows the same rules.}
{An \cmd{ULTIMATE\_*} result.}
{Adds an array and enters it. Keys on items inside the array are ignored.}
{static const char key[] = \{0x73,0x63,0x6f,0x72,0x65,0x73,0\};\
err = ultimate\_http\_body\_array(body, key);}
{Following typed values become array elements until a body-up call.}

\capiref{ultimate\_http\_body\_binary}{ultimate.h}
{uint8\_t ultimate\_http\_body\_binary(uint8\_t handle, const uint8\_t *data, uint16\_t len)}
{\cmd{handle} names an \cmd{HTTP\_BODY\_BINARY} slot; \cmd{data} supplies
\cmd{len} bytes.}
{An \cmd{ULTIMATE\_*} result.}
{Appends raw bytes; successive calls extend the same body.}
{err = ultimate\_http\_body\_binary(body, chunk, sizeof chunk);}
{The chunk is appended byte-for-byte.}

\begin{tip}
Build a large JSON document a field at a time inside the Ultimate, exchange it,
then free both body and request handles on every path. This avoids keeping the
whole document in the C64's limited RAM.
\end{tip}

\section{Errors And Diagnostics}

\capirefout{ultimate\_strerror}{ultimate.h}
{const char *ultimate\_strerror(uint8\_t err)}
{\cmd{err} is an \cmd{ULTIMATE\_*} result.}
{A non-null pointer to stable PETSCII English text.}
{Formats SDK-level diagnostics without exposing target-specific meanings.}
{printf("load: \%s\textbackslash n", ultimate\_strerror(ULTIMATE\_ERR\_IO));}
{LOAD: I/O ERROR}

\cmacroref{ultimate\_device\_code}{ultimate\_device\_code()}
{None.}
{The most recent raw target status, or \cmd{UCI\_DEVICE\_CODE\_NONE}.}
{Aliases \cmd{uci\_last\_device\_code()} for diagnostics.}
{uint16\_t device = ultimate\_device\_code();}
{\cmd{device} holds firmware-specific detail; application control should use
the SDK result instead.}

\section{Low-Level UCI Transport}

\capiref{uci\_signature\_present}{uci.h}
{uint8\_t uci\_signature\_present(void)}
{None.}
{1 when the documented identification signature is present, otherwise 0.}
{Performs a quick register-only probe. It cannot distinguish a healthy from a
wedged interface.}
{if (!uci\_signature\_present()) return 1;}
{The program exits when no UCI signature is mapped.}

\capiref{uci\_ident\_register}{uci.h}
{uint8\_t uci\_ident\_register(void)}
{None.}
{The raw identification-register byte.}
{Exposes register detail for diagnostics and new bindings.}
{uint8\_t ident = uci\_ident\_register();}
{\cmd{ident} contains the unprocessed hardware value.}

\capiref{uci\_init}{uci.h}
{uint8\_t uci\_init(void)}
{None.}
{\cmd{ULTIMATE\_OK}, \cmd{ULTIMATE\_ERR\_NO\_DEVICE}, or
\cmd{ULTIMATE\_ERR\_TIMEOUT}.}
{Checks the signature, clears latent state, aborts an abandoned exchange if
needed, and installs the default timeout.}
{uint8\_t err = uci\_init();}
{The raw transport is idle and ready only after success.}

\capiref{uci\_exec}{uci.h}
{uint8\_t uci\_exec(uci\_request *req)}
{\cmd{req} points to a fully initialized
\hyperref[struct:uci-request]{\cmd{uci\_request}} containing target, command,
input spans and reply buffers.}
{An \cmd{ULTIMATE\_*} result. \cmd{ULTIMATE\_ERR\_TRUNCATED} means the command
completed but a non-null data buffer filled; output lengths are always updated.}
{Runs one command through every reply block. With
\cmd{UCI\_TARGET\_NO\_REPLY}, waits only for idle and collects nothing.}
{err = uci\_exec(\&req);}
{\cmd{req.datalen} and \cmd{req.statuslen} describe the retained reply.}

\capiref{uci\_exec\_first}{uci.h}
{uint8\_t uci\_exec\_first(uci\_request *req)}
{\cmd{req} is the same
\hyperref[struct:uci-request]{\cmd{uci\_request}} layout used by
\cmd{uci\_exec()}.}
{An \cmd{ULTIMATE\_*} result for the first reply block; \cmd{req->datalen} is
reset to this block's length.}
{Starts a boundary-preserving exchange such as directory reading.}
{err = uci\_exec\_first(\&req);}
{The first logical block is available without being joined to the next.}

\capiref{uci\_exec\_next}{uci.h}
{uint8\_t uci\_exec\_next(uci\_request *req)}
{\cmd{req} is the same
\hyperref[struct:uci-request]{\cmd{uci\_request}} passed to
\cmd{uci\_exec\_first()}; no other command may intervene.}
{An \cmd{ULTIMATE\_*} result for the next block; final status is decoded on the
last.}
{Advances one live boundary-preserving exchange.}
{if (uci\_more\_blocks()) err = uci\_exec\_next(\&req);}
{The next block replaces the previous contents and length.}

\capiref{uci\_more\_blocks}{uci.h}
{uint8\_t uci\_more\_blocks(void)}
{Call after \cmd{uci\_exec\_first()} or \cmd{uci\_exec\_next()}.}
{1 when another block follows, otherwise 0.}
{Reads current transport state without issuing a command.}
{while (uci\_more\_blocks()) err = uci\_exec\_next(\&req);}
{The loop consumes all remaining blocks.}

\capiref{uci\_abort}{uci.h}
{uint8\_t uci\_abort(void)}
{None.}
{An \cmd{ULTIMATE\_*} result.}
{Abandons an exchange and forces the interface back to idle. Safe at any time.}
{err = uci\_abort();}
{A half-read directory or stopped program no longer holds the interface.}

\capiref{uci\_set\_timeout}{uci.h}
{void uci\_set\_timeout(uint8\_t units)}
{\cmd{units} counts groups of 256 status polls; zero means forever.}
{Nothing.}
{Sets the polling budget. Its wall time changes with CPU speed.}
{uci\_set\_timeout(UCI\_TIMEOUT\_DEFAULT);\
uci\_set\_timeout(UCI\_TIMEOUT\_FOREVER);}
{The calls demonstrate bounded default and unbounded operation.}

\capiref{uci\_get\_timeout}{uci.h}
{uint8\_t uci\_get\_timeout(void)}
{None.}
{The current timeout units; zero means forever.}
{Reads the transport's current polling budget.}
{uint8\_t saved = uci\_get\_timeout();}
{\cmd{saved} can restore the caller's policy after a temporary change.}

\capiref{uci\_last\_device\_code}{uci.h}
{uint16\_t uci\_last\_device\_code(void)}
{None.}
{The last decoded raw status: DOS-family 0--99, HTTP 0--999, SoftwareIEC
0--255, or \cmd{UCI\_DEVICE\_CODE\_NONE}.}
{Exposes target-specific diagnostic detail.}
{uint16\_t device = uci\_last\_device\_code();}
{Interpret \cmd{device} only in the context of the target just called.}

\capiref{uci\_status\_format}{uci.h}
{uint8\_t uci\_status\_format(uint8\_t target)}
{\cmd{target} is a \cmd{UCI\_TARGET\_*} value.}
{The expected \cmd{UCI\_STATUS\_FMT\_*} encoding.}
{Reports protocol documentation; it does not inspect bytes and must not be used
as the decoder because some firmware replies violate the target convention.}
{uint8\_t dos = uci\_status\_format(UCI\_TARGET\_DOS1);\
uint8\_t http = uci\_status\_format(UCI\_TARGET\_HTTP);\
uint8\_t iec = uci\_status\_format(UCI\_TARGET\_SOFTIEC);}
{The values are \cmd{UCI\_STATUS\_FMT\_DECIMAL2},
\cmd{UCI\_STATUS\_FMT\_DECIMAL3}, and \cmd{UCI\_STATUS\_FMT\_BINARY},
respectively. Other target numbers use the two-digit default.}

\capiref{uci\_decode\_status}{uci.h}
{uint8\_t uci\_decode\_status(uint8\_t target, const uint8\_t *status, uint16\_t statuslen)}
{\cmd{target} supplies context; \cmd{status} points to \cmd{statuslen} raw bytes.
Pass the complete status when possible; a nonzero length requires a non-null
pointer.}
{The decoded \cmd{ULTIMATE\_*} result and, as a side effect, a new raw device
code. Empty status is success.}
{Detects two- and three-digit numeric forms, HTTP response lines, SoftwareIEC
binary status, and unnumbered target-error text from the bytes rather than
trusting the expected target format.}
{static const uint8\_t dos[] = \{0x30,0x30,0x2c,0x4f,0x4b\};\
static const uint8\_t http[] = \{0x34,0x30,0x34,0x20\};\
static const uint8\_t line[] =\
\hspace*{2em}\{0x48,0x54,0x54,0x50,0x2f,0x31,0x2e,0x31,0x20,0x32,0x30,0x30\};\
static const uint8\_t iec[] = \{0x00\};\
static const uint8\_t text[] = \{0x46,0x55,0x4c,0x4c\};\
err = uci\_decode\_status(UCI\_TARGET\_DOS1, dos, sizeof dos);\
err = uci\_decode\_status(UCI\_TARGET\_HTTP, http, sizeof http);\
err = uci\_decode\_status(UCI\_TARGET\_HTTP, line, sizeof line);\
err = uci\_decode\_status(UCI\_TARGET\_SOFTIEC, iec, sizeof iec);\
err = uci\_decode\_status(UCI\_TARGET\_DOS1, text, sizeof text);\
err = uci\_decode\_status(UCI\_TARGET\_DOS1, NULL, 0);}
{The calls demonstrate two-digit status, three-digit firmware HTTP status, an
HTTP response line, binary status, unnumbered target-error text, and empty
success. \cmd{uci\_last\_device\_code()} retains a decoded number or
\cmd{UCI\_DEVICE\_CODE\_NONE} when none exists.}

\evidence{\cmd{include/ultimate.h}; \cmd{include/uci.h};
\cmd{guide/examples/c\_calls.c}; \cmd{tests/emulator/sdk.suite};
\cmd{tests/hardware/ucitest.c}}

\end{referencechapter}
